% Source PDF: source/普通高考/2014/2014辽宁理.pdf
% Page: 1 (program flowchart for question 13).
% PDF embedded bitmap attempt: pdfimages -list found no embedded image; the source figure is vector content on the page.
% Grid evidence: tmp/redraw_flowchart_2014_liaoning_science/recheck_source/page-1-grid100.png (250 dpi; 100 px major grid used for locating).
% No-grid crop evidence: tmp/redraw_flowchart_2014_liaoning_science/recheck_source/q13_orig_tight.png.
% Elements (complete): rounded START and END terminals; INPUT x and OUTPUT y parallelograms;
% process boxes y=x/3+2 and x=y; decision diamond |y-x|<1; directed connectors and labels 是/否.
% Node -> edge checklist (complete): START -> INPUT x -> y=x/3+2 -> decision;
% decision(是) -> OUTPUT y -> END; decision(否) -> x=y -> y=x/3+2.
% All connectors are solid directed edges; no implicit, duplicate, or decorative arrows.

\begin{tikzpicture}[
    baseline=-0.5ex,
    >=Stealth,
    line width=0.45pt,
    line cap=round,
    line join=round,
    font=\normalsize,
    flowbox/.style={draw, minimum height=0.56cm, inner sep=1.5pt, align=center,
        execute at begin node={\setlength{\parindent}{0pt}}},
    terminal/.style={flowbox, rounded corners=3pt, minimum width=1.22cm},
    io/.style={flowbox, trapezium, trapezium left angle=72, trapezium right angle=108,
        minimum width=1.70cm},
    process/.style={flowbox, minimum width=2.55cm},
    decision/.style={flowbox, diamond, aspect=2.75, inner sep=0.6pt,
        minimum width=3.45cm, minimum height=0.78cm},
]
    \node[terminal] (start) at (0,4.45) {开始};
    \node[io]       (input) at (0,3.55) {输入 $x$};
    \node[process]  (calc)  at (0,2.45) {$y=\dfrac{x}{3}+2$};
    \node[decision] (test)  at (0,1.05) {$|y-x|<1$};
    \node[io]       (output)at (0,-0.30) {输出 $y$};
    \node[terminal] (finish)at (0,-1.22) {结束};
    \node[process]  (assign)at (2.95,2.45) {$x=y$};

    \draw[->] (start.south) -- (input.north);
    \draw[->] (input.south) -- (calc.north);
    \draw[->] (calc.south) -- (test.north);
    \draw[->] (test.south) -- node[right=1pt] {是} (output.north);
    \draw[->] (output.south) -- (finish.north);
    \draw[->] (test.east) -- node[below right=1pt] {否} (2.95,1.05) -- (assign.south);
% The return edge joins the existing INPUT-to-calc trunk without a duplicate arrowhead.
    \draw (assign.north) -- ++(0,0.25) -- (0,2.98);
\end{tikzpicture}
